\documentclass[11pt]{article}
\usepackage[T1]{fontenc}\usepackage{lmodern,amsmath,amssymb,amsthm}
\usepackage[margin=1in]{geometry}\usepackage[hidelinks]{hyperref}
\newcommand{\Q}{\mathbb Q}\newcommand{\Z}{\mathbb Z}\newcommand{\F}{\mathbb F}
\newcommand{\A}{\mathbb A}\newcommand{\Sha}{\mathop{\mathrm{Sha}}\nolimits}
\DeclareMathOperator{\Sel}{Sel}\DeclareMathOperator{\Br}{Br}
\newtheorem{theorem}{Theorem}[section]\newtheorem{lemma}[theorem]{Lemma}\newtheorem{corollary}[theorem]{Corollary}
\title{A recurrence test for the Brauer--Manin sets of three quartic coverings}
\author{\small\href{mailto:admin@cosmosgeek.org}{admin@cosmosgeek.org}\quad
\href{mailto:b321014@iiit-bh.ac.in}{b321014@iiit-bh.ac.in}}
\date{Working paper for specialist review}
\begin{document}\maketitle
\begin{abstract}
Let $p\equiv1\pmod 8$ be prime. Three locally soluble isogeny coverings of $y^2=x^3-4px$ are torsors under $E:y^2=x^3+px$. We show that the full Brauer--Manin set of each is nonempty exactly when Nomoto's coefficient $f_{3(p-1)/8}(0)$ vanishes modulo $p$. The same test separates two possible structures of the two-primary Selmer groups, although the first two-descent has dimension three in both cases. The proof combines Nomoto's congruence, the rank-zero CM converse at $2$, the Cassels--Tate pairing and Creutz's divisibility criterion. We also express the two-primary valuation of a rank-zero CM central value as a sum of the dimensions of the successive Selmer images. Its minimal valuation is read by a single Cassels--Tate Pfaffian.
\end{abstract}

\section{Statement}
Fix a prime $p\equiv1\bmod8$, and define
\[
 E:y^2=x^3+px,\qquad E':y^2=x^3-4px,\qquad r=\operatorname{rank}E(\Q).
\]
Write $C_d$ for the smooth projective model of $z^2=du^4-(4p/d)v^4$ in weighted projective coordinates $(u:v:z)$ of weights $(1,1,2)$. Define the integer polynomials $f_n$ by
\begin{align*}
 f_0(t)&=1,\qquad f_1(t)=2t+3,\\
 f_{n+1}(t)&=(4n+1)(2t+3)f_n(t)-12(t+1)(t+2)f'_n(t)\\
 &\qquad -2n(2n-1)(t^2+3t+3)f_{n-1}(t)\qquad(n\ge1),
\end{align*}
and set $F_p=f_{3(p-1)/8}(0)$. For a curve $C$ the set $C(\A_{\Q})^{\Br}$ consists of adelic points orthogonal to its entire cohomological Brauer group under the Brauer--Manin pairing.

\begin{theorem}\label{main}
For each $d\in\{-1,2,-2\}$ one has $C_d(\A_{\Q})\ne\varnothing$ and
\[
 C_d(\A_{\Q})^{\Br}\ne\varnothing
 \quad\Longleftrightarrow\quad L(E,1)=0
 \quad\Longleftrightarrow\quad p\mid F_p .
\]
Exactly $2^r-1$ of these three coverings have a rational point. If $p\nmid F_p$ then $r=0$. If $p\mid F_p$ then the analytic rank is even and at least two, and $r\le2$.
\end{theorem}

\section{From the recurrence to a central value}
For the primitive positive real period $\Omega_1$ of $y^2=x^3+x$, put $S_p=2p^{1/4}L(E,1)/\Omega_1$. Nomoto's dissertation \cite[\S5.2]{thesis} gives $S_p\in\Z$ by Birch and Swinnerton-Dyer's integrality theorem. Its Proposition~5.8 and Remark~5.9 identify $S_p$ modulo a prime above $p$, up to a unit, with the algebraic Hecke central value $L_{E,k}$ for $k=(3p+1)/4$. The factorial in Remark~5.9 is a unit at $p$. Theorem~7.6 identifies $L_{E,k}$ with $f_{(k-1)/2}(0)^2$ because $k$ is odd. We use the recurrence as stated in \cite[Thm.~1.1]{nom} and \cite[Thm.~4.2]{thesis}. Thus $p\mid S_p$ is equivalent to $p\mid F_p$.

The size bound $|S_p|<p$ in \cite[Prop.~5.6]{thesis} turns the congruence into an exact zero test. We give an elementary estimate for this family. Its conductor is $64p^2$ and its root number is $+1$ \cite[\S4--5]{thesis}. At the centre the weight-two functional equation gives the exact series
\[
 L(E,1)=2\sum_{n\ge1}\frac{a_n}{n}\exp(-\pi n/(4p)).
\]
Indeed splitting the Mellin integral at $1/\sqrt{64p^2}$ and applying the Fricke involution makes the two integrals equal. In the CM ideal sum for $a_n$, each contributing ideal of norm $n$ prime to $2p$ has coefficient of absolute value $\sqrt n$ \cite[\S5.2]{thesis}. In $\Z[i]$ its four generators have opposite coordinate parities. At $\max(|a|,|b|)=j\ge1$ there are $4j$ pairs $(a,b)$ of opposite parities, each of norm at least $j^2$. Hence
\[
 |L(E,1)|\le 2\sum_{j\ge1}\exp(-\pi j^2/(4p))<2\sqrt p .
\]
The real period $\Omega_1=\int_0^\infty (x^3+x)^{-1/2}dx$ is at least $2\sqrt2$ by integrating separately on $[0,1]$ and $[1,\infty)$. Since $p\ge17$, $4p^{3/4}/\Omega_1<p$ follows. In particular
\begin{equation}\label{test} L(E,1)=0\quad\Longleftrightarrow\quad p\mid F_p .\end{equation}

\section{The two isogenies and their local conditions}
Let $\phi:E\to E'$ have kernel generated by $(0,0)$ and let $\widehat\phi$ be its dual. For an isogeny $\psi:A\to B$, we write $\Sel_\psi(A)$ for the subgroup of $H^1(\Q,A[\psi])$ cut out by the local images of $B/\psi A$. For these rational two-isogenies the connecting maps are represented by the abscissa on ordinary affine points. At $(0,0)$ they take the squareclass of the coefficient of $x$ in the equation of the \emph{target} \cite[X.4.9]{sil}. In particular
\[
 \delta_\phi(E'[\widehat\phi](\Q))=\langle-p\rangle,\qquad
 \delta_{\widehat\phi}(E[\phi](\Q))=\langle p\rangle.
\]

The $\phi$-cover with parameter $d$ is $C_d$. Over $v\ne0$, its map to $E'$ is
\[
 (u:v:z)\longmapsto
 \bigl(d u^2/v^2,\ d u z/v^3\bigr).
\]
Indeed its pullback adjoins a square root of $x'/d$. The original isogeny adjoins a square root of $x'$, as $x'=(y/x)^2$ under $\phi$. The deck group is $E[\phi]$, acting on $E$ by translation. The class of the locally soluble $C_d$ in the Weil--Ch\^atelet group therefore lies in $\Sha(E)[\phi]$.
The dual coverings instead have equation $z^2=d u^4+(p/d)v^4$ and map to $E$.

\begin{lemma}\label{local}
In $\Q^\times/\Q^{\times2}$ the two isogeny Selmer groups are
\[
 \Sel_\phi(E)=\langle-1,2,-p\rangle,\qquad
 \Sel_{\widehat\phi}(E')=\langle p\rangle .
\]
\end{lemma}
\begin{proof}
At any prime $q\nmid 2p$, a class of odd $q$-adic valuation is excluded: the valuations of the two terms of its biquadratic differ modulo four by two, so their minimum is odd. Here we rescale to primitive local coordinates; zero coordinates give the same exclusion immediately. This bounds the $\phi$-group by $\langle-1,2,p\rangle$. For the dual coverings the same calculation applies at $2$, since $p$ is odd. The real place requires positive $d$ since the right-hand side is negative if $d<0$. This gives the dual upper bound, attained by its torsion class.

For the $\phi$-group the class $-p$ is represented by the rational point $(0,0)\in E'(\Q)$. We verify $-1$ and $2$. At $p$, choose $v=0$ since both are squares in $\Q_p$. At infinity choose $u=0$ for $d=-1$ and $v=0$ for $d=2$. At $2$, for $d=-1$ take $u=0,v=1$, as $p\equiv1\bmod8$ is a $2$-adic square. For $d=2$, put $v=1$. If $p\equiv1\bmod16$ take $u^4=p,z=0$, and if $p\equiv9\bmod16$ take $u^4=p+8,z=4$. We used $(\Z_2^\times)^4=1+16\Z_2$. At every other prime the binary quartic is separable in odd residue characteristic and defines a smooth genus-one curve with a residue field point by the Hasse bound. This point lifts. The local Kummer images are groups, so their conditions also hold on the product class $-2$.
\end{proof}

\begin{lemma}\label{dim}
The group $\Sel_2(E)$ has $\F_2$-dimension three. The group $E(\Q)[2^\infty]$ has order two, and $\Sel_{2^\infty}(E)[2]$ has dimension two.
\end{lemma}
\begin{proof}
The isogeny Selmer sequence is exact at its first three nonzero terms:
\[
 0\longrightarrow E'[\widehat\phi](\Q)/\phi E[2](\Q)
 \longrightarrow\Sel_\phi(E)\longrightarrow\Sel_2(E)
 \longrightarrow\Sel_{\widehat\phi}(E').
\]
The initial quotient has order two. The last map surjects here since the dual Selmer group is generated by the image of the rational point $(0,0)$ on $E$, which lifts via its $2$-Kummer class. Lemma~\ref{local} now gives dimension $3-1+1=3$.
The rational two-torsion of $E$ consists of $O$ and $(0,0)$. For a half of $(0,0)$ the duplication formula would give $(x^2-p)^2=0$, impossible over $\Q$. The usual Kummer comparison
\[
 0\longrightarrow E(\Q)[2^\infty]/2E(\Q)[2^\infty]
 \longrightarrow \Sel_2(E)\longrightarrow\Sel_{2^\infty}(E)[2]
 \longrightarrow0
\]
gives the last assertion.
\end{proof}

\section{The full Brauer--Manin obstruction}
Write $S=\Sel_{2^\infty}(E)$, $U=\Q_2/\Z_2$ and $T=\Sha(E)[2^\infty]$. The exact sequence
\[
 0\longrightarrow E(\Q)\otimes U\longrightarrow S\longrightarrow T\longrightarrow0
\]
splits as a sequence of abelian groups since its first term is divisible. Thus $r\le2$. The cofinite group $S$ is a sum of copies of $U$ and a finite group. The latter is isomorphic to $T/T_{\rm div}$. It has a perfect alternating Cassels--Tate pairing for an elliptic curve \cite{ps}, so its cyclic elementary divisors occur in equal pairs. Lemma~\ref{dim} gives
\[
 \dim_{\F_2} S[2]=2
   =\operatorname{corank}_{\Z_2}S+\dim_{\F_2}(T/T_{\rm div})[2].
\]
In particular the corank is zero or two.

If $L(E,1)=0$, the rank-zero CM converse of Burungale--Tian \cite[Thm.~1.1]{bt} rules out corank zero. That theorem applies at every prime to rational CM curves. Here the CM field is $\Q(i)$ and the Selmer prime is $2$, including the additive reduction at that prime. It follows that $S\simeq U^2$ and $T\simeq U^{2-r}$. Every class represented by $C_d$ is divisible in $\Sha(E)$. By Creutz \cite[Prop.~18 and proof of Thm.~5]{creutz}, for a locally soluble torsor under an abelian variety, its full Brauer--Manin set is nonempty exactly when its class belongs to the maximal divisible subgroup of $\Sha$. Hence each of the three sets is nonempty.

If $L(E,1)\ne0$, the analytic-rank-zero theorem gives $r=0$ and finite $\Sha(E)$ \cite[\S1]{bt}. The curve $E'$ then also has rank zero. Both curves have just one nonzero rational point of order two, and neither has a point of order four: halving $(0,0)$ would require $x^2=p$ or $x^2=-4p$, respectively. Odd torsion contributes nothing to the isogeny quotient. Consequently the image of $E'(\Q)/\phi E(\Q)$ is exactly $\langle-p\rangle$. Each $d=-1,2,-2$ represents a distinct nonzero class in $\Sha(E)[\phi]$, and no such class is divisible in finite $\Sha(E)$. Creutz's criterion makes each Brauer--Manin set empty.

For either branch the image of $E(\Q)/\widehat\phi E'(\Q)$ is exactly $\langle p\rangle$, by Lemma~\ref{local} and rational two-torsion. The analogous sequence of rational point quotients gives
\[
 0\longrightarrow E'[\widehat\phi](\Q)/\phi E[2](\Q)
 \longrightarrow E'(\Q)/\phi E(\Q)\longrightarrow E(\Q)/2E(\Q)
 \longrightarrow E(\Q)/\widehat\phi E'(\Q)\longrightarrow0 .
\]
Thus $|E'(\Q)/\phi E(\Q)|=2^{r+1}$. Modulo its torsion class $-p$, this image has $2^r$ elements, and $1,-1,2,-2$ represent all the ambient $\phi$-Selmer cosets modulo $\langle-p\rangle$. This proves the rational-point count and completes Theorem~\ref{main}.

\begin{corollary}
For every $n\ge1$ the two branches of the recurrence yield
\[
\begin{array}{c|c|c}
 & \Sha(E)[2^\infty] & |\Sel_{2^n}(E)|\\ \hline
 p\nmid F_p & (\Z/2^e\Z)^2\ (e\ge1) & 2^{1+2\min(e,n)}\\
 p\mid F_p & U^{\,2-r} &2^{1+2n}.
\end{array}
\]
\end{corollary}
\begin{proof}
On the nonvanishing branch the finite two-primary group has two generators by Lemma~\ref{dim} and rank zero. Pairing duality forces equal exponents. On both branches the finite Kummer sequence gives
\[
 |\Sel_{2^n}(E)|=2^{nr+1}|\Sha(E)[2^n]| ,
\]
since $|E(\Q)/2^nE(\Q)|=2^{nr+1}$.
\end{proof}

\section{A Pfaffian test for the normalized central value}

The following calculation applies to a CM elliptic curve $A/\Q$ with geometric endomorphism ring the full ring of integers of its CM field.
Write $T=A(\Q)_{\rm tors}$ and
\[
 V=\Sel_2(A/\Q)/\operatorname{im}(T),\qquad
 R_j=\operatorname{im}\bigl(\Sel_{2^j}(A/\Q)\longrightarrow V\bigr)
 \quad(j\ge1).
\]
The transition is multiplication by $2^{j-1}$ on coefficient modules
followed by projection. The torsion image uses the entire group $T$
under $A(\Q)/2A(\Q)\to\Sel_2$. This convention covers rational
points of order four. Put $n=\dim_{\F_2}V$, so $R_1=V$, and define
\[
 J(A)=\frac{|T|^2L(A,1)}{\Omega_A\prod_{p<\infty}c_p(A)},\qquad
 \Omega_A=\int_{A(\mathbb R)}|\omega_{\min}|.
\]
The integral includes all real components. In an $\F_2$-basis of $V$
let $C$ be the alternating matrix of the Cassels--Tate pairing pulled back
from $\Sha(A/\Q)[2]$, identifying $\tfrac12\Z/\Z$ with $\F_2$.

\begin{theorem}\label{th:cm-pfaffian}
With $v_2(0)=+\infty$ one has
\[
 v_2(J(A))=\sum_{j\ge1}\dim_{\F_2}R_j.
\]
If $n$ is even, $2^{-n}J(A)$ is an integer and
\[
                 2^{-n}J(A)\equiv\operatorname{Pf}(C)\pmod 2.
\]
If $n=0$ the empty Pfaffian is $1$.
\end{theorem}

\begin{proof}
Put $X=\Sel_{2^\infty}(A/\Q)$ and write
$X\simeq(\Q_2/\Z_2)^\rho\oplus\bigoplus_i\Z/2^{e_i}\Z$
as an abstract group, with $e_i\ge1$. The exact coefficient sequence
gives a surjection $\Sel_{2^j}\to X[2^j]$ with kernel the image
of $A(\Q)[2^\infty]/2^j A(\Q)[2^\infty]$; surjectivity follows
also by comparing the finite and infinite Kummer sequences.
Consequently $V=X[2]$ and $R_j=2^{j-1}X[2^j]$. Thus
$\dim R_j=\rho+\#\{i:e_i\ge j\}$. The divisible Mordell--Weil
subgroup of $X$ splits as an abstract direct summand. The remaining
finite elementary divisors are those of $\Sha[2^\infty]$ modulo its
divisible subgroup. The latter finite quotient has a perfect
alternating Cassels--Tate pairing, hence its elementary divisors occur
in equal pairs. Restricted to its subgroup killed by two, its radical
is the intersection with twice the group. The divisible part and the
Mordell--Weil part pair to zero. It follows that
$\operatorname{rad}(C)=R_2$.

In Burungale--Flach take the ground field $F=K$ equal to the imaginary quadratic CM field and $F^+=\Q$. All endomorphisms are defined over $K$ and its torsion action is abelian by complex multiplication. If $L(A,1)\ne0$, Corollary~2 gives $A(\Q)$ and
$\Sha(A/\Q)$ finite and $J(A)=|\Sha(A/\Q)|$ \cite{pfBF}. Then
$\rho=0$ and both sides of the valuation identity equal $\sum_i e_i$.
If $L(A,1)=0$, the CM corank-zero converse of Burungale--Tian
\cite{bt} at $p=2$ implies $\rho>0$, so both sides are infinite.
For nonzero $J$, $v_2(J)=n$ precisely when $R_2=0$. For zero $J$,
$R_2$ contains the two-torsion of the nonzero divisible part. Thus
the normalized integer is odd precisely when $C$ is nonsingular.
Over $\F_2$ this is equivalent to $\operatorname{Pf}(C)=1$.
\end{proof}

In a marked family for which the matrices are independently known to be
$C+B$, the theorem evaluates $\operatorname{Pf}(C+B)$ pointwise.
As a polynomial in independent upper-triangular variables $B_{ij}$, its
lowest nonzero homogeneous degree is $\dim\operatorname{rad}(C)/2$.
Indeed, put $C$ in symplectic normal form. Every matching monomial must
supply variable entries pairing all radical indices, and the bound is
attained by pairing the radical indices among themselves and taking the
constant symplectic pairs on the complement. This order is invariant
under a change of basis. In the original basis the coefficient of a
fixed matching on the complement of $\{i,j\}$ is $C_{ij}$. On the
two-dimensional reduced space of the quartic family this reduces to
one pairing bit. In higher dimensions the lowest degree is at least
half the rational rank, since all Mordell--Weil classes lie in the
radical.

This multiplicity can be arbitrarily large at an analytically
nonvanishing twist of one fixed CM curve. Take $A_0:y^2=x^3-2$.
It has $j=0$ and no rational point of order two; its two-division field
has group $S_3$. For any $h\ge1$ and $j_0\ge2$, Smith's Theorem~1.5
\cite{pfSmith} gives a positive density of integer twist parameters
with ranks $r_2=r_4=\cdots=r_{2^{j_0}}=2h$ and
$r_{2^{j_0+1}}=0$. Indeed the initial probability for $2h$ is
positive, a zero alternating matrix preserves that state, and an
invertible alternating matrix takes it to zero. Here $r_{2^j}=\dim R_j$
by the finite coefficient description above and Smith's definition.
Therefore $C=0$, $\rho=0$, and $L(A_0^d,1)\ne0$. At this matrix the
Pfaffian has multiplicity $h$, and the exact two-adic valuation of
$J(A_0^d)$ is $2hj_0$. Thus arbitrarily many degenerate layers are
compatible with analytic rank zero.

For comparison, suppose $A[2]$ has full $S_3$ image. Morgan's
Section~2 sets up an arbitrary principally polarized abelian variety;
Section~5.2 requires (A) and (B) of his Assumption~1.3
\cite{pfMorgan}. These hold for the natural $S_3$ module. Also
$H^1(S_3,A[2])=0$ by inflation--restriction through $C_3$,
and the symmetric bundle $\mathcal O_A(O)$ gives $\mathfrak c=0$.
Every alternating pairing on the two-Selmer space is admissible:
$\dim V-\rho$ is even by the paired elementary divisors just described.
Theorem~5.8 thus supplies quadratic twists with a prescribed matrix on
the same marked Selmer group. Over $\Q$, a positive-density set of positive prime twists realizes each specified alternating matrix $P$ on this marked space. To see this, let $F$ be the governing field in the proof of Morgan's Theorem~5.8, with $T$ the chosen basis and $T'$ empty. Let $\Sigma$ contain $2$, all bad primes and the primes ramified in $F$. Take the compositum $U$ of $\Q(\zeta_8)$ and the unique quadratic subfields of $\Q(\zeta_l)$ for odd $l\in\Sigma$. This is an elementary abelian $2$-extension. Morgan's Proposition~4.14 puts all quadratic subfields of $F$ in $\Q(A[2])$. Thus $F\cap U$ is contained in the unique quadratic subfield of the $S_3$-extension. His Corollary~4.19 gives a lift $\sigma$ of a three-cycle with class-function values $\psi_{ij}(\sigma)=P_{ij}-C_{ij}$ for all $i<j$. Here the numbering is that of the published paper, corresponding to Proposition~4.13 and Corollary~4.18 of the preprint.

This lift fixes $F\cap U$, so $(\sigma,1)$ defines an element of $\operatorname{Gal}(FU/\Q)$. Its conjugacy class has positive Chebotarev density. Each corresponding prime $q\notin\Sigma$ is $1\bmod 8$ and a square in $\Q_l$ at all odd $l\in\Sigma$. Its quadratic twist character is trivial at $\Sigma$ and infinity, and ramifies only at $q$, where Frobenius fixes no nonzero vector of $A[2]$. Morgan's Lemma~2.8 preserves the local Selmer conditions, and Proposition~3.3 supplies the indicated differences of the pairings. This gives the claimed prime twists over $\Q$.

Finally, suppose $A$ has good supersingular reduction at $2$,
$A(\Q_2)[2]=0$, and its matrix has radical a line spanned by $\beta$.
If $\beta$ has nonzero local Kummer image at $2$, it is represented by
a rational point spanning $A(\Q)/2A(\Q)$, and $A$ has analytic rank
one. Indeed $\dim R_2=1$ forces $\rho=1$ because the finite elementary
divisors occur in pairs. Thus $R_2$ is the two-torsion in the divisible
line of $X$. The local torsion hypothesis identifies
$A(\Q_2)/2A(\Q_2)$ with the two-torsion of
$A(\Q_2)\otimes\Q_2/\Z_2$. Nonzero localization on that line implies
surjectivity on the divisible group. Burungale--Skinner's Appendix~A,
Theorem~A.1 and Remark~A.5 \cite{pfBS}, apply also at $p=2$, giving
analytic and rational rank one. Here the CM and good-reduction hypotheses are essential. The Mordell--Weil line, being contained
in the radical, is then exactly the line spanned by $\beta$.

One can first lower an even dimension $n\ge2$ by one while preserving a specified class $\beta$. Suppose in addition that complex conjugation acts as a transposition on $A[2]$, the curve has good supersingular reduction at two, $A(\Q_2)[2]=0$, and the localization of $\beta$ at two is nonzero. In this paragraph $A[2]$ has full $S_3$ image, as in the preceding prime-twist construction. Choose a nonzero linear functional on $\Sel_2(A)$ killing $\beta$. Mazur--Rubin's evaluation lemma \cite[Lem.~3.5]{mr}, with complex conjugation as the reference element, realizes that functional as the localization at a good prime $q$ with one-dimensional local Kummer image. More explicitly, choose the corresponding Frobenius on the field cut out by the Selmer cocycles, with cyclotomic action equal to complex conjugation on $\Q(\zeta_M)$, where $M=8\prod_{l\mid N_A,\,l\ {\rm odd}} l$. This simultaneous choice follows from the lemma with $\Q(A[2])\Q^{\rm ab}$ fixed. Here $N_A$ is the conductor.

Now $q\equiv-1\bmod M$ and the twist by $-q$ is locally trivial at every prime dividing $2N_A$. The real local cohomology vanishes for the nontrivial transposition module over $\F_2$. At $q$ the transverse Kummer line replaces the unramified line, and the old localization is surjective. The local comparison in \cite[Prop.~3.3]{mr} shows that the new Selmer space is precisely the kernel of the chosen functional. Its odd dimension is $n-1$. An alternating matrix on that kernel can have radical exactly $\langle\beta\rangle$. Applying the positive prime-twist construction to $A^{-q}$ and keeping $2N_Aq$ locally split realizes this radical for infinitely many twists $A^{-qq'}$. The localization at two remains nonzero, so these twists have analytic and rational rank one. Each quadratic field $\Q(\sqrt{-qq'})$ is imaginary with all primes dividing $N_A$ split. The special family of Theorem~\ref{main} has rational two-torsion; this additional construction concerns the full-$S_3$ case.

For these rank-one twists let $K=\Q(\sqrt d)$, write $c$ for its involution and identify $A^d(\Q)$ with $B^-$ in $B=A(K)$; take $B^+=A(\Q)$. Norm and difference inject $B/(B^++B^-)$ into $B^+/2B^+$ and $B^-/2B^-$ respectively: for example $(1+c)T=2P$ implies $T-P\in B^-$. Thus the quotient has order at most two. Any nonzero class gives $P=(1+c)T$ and $R=(1-c)T$ whose Kummer classes agree over $\Q$. Indeed for a half $H$ of $P$, $T-H$ is a twisted half of $R$ and the two cocycles agree in $A[2]$. Conversely if $\beta$ is represented by $P$ on $A$ and a generator $R$ modulo torsion on the twist, restriction to $K$ gives $P+R=2T$ there. No quadratic field contains a nonzero point of $A[2]$ since its Galois action is $S_3$. Hence $(1+c)T=P$ and $(1-c)T=R$ exactly. It follows that the quotient has order two if $\beta$ comes from $A(\Q)$ and one otherwise. For $L(A,1)=L'(A,1)=0$ the product formula for quadratic base change gives $L'''(A/K,1)=3L''(A,1)L'(A^d,1)$ with nonzero last factor. This identity leaves the base second derivative undetermined.

\section{A radial moment at fixed weight}
The recurrence reads a central \emph{value}. For comparison, the complex Mellin derivatives can be expressed using an additional index in a theta expansion. Let $\Lambda\subset\mathbb C$ be an ideal lattice and $w$ a periodic function on it encoding the finite part of a Hecke character of infinity type one. Fix $c>0$, put $X_\alpha=c|\alpha|^2$ and take the cuspidal theta kernel
\[
 \Theta(\rho)=\sum_{\alpha\in\Lambda}w(\alpha)\alpha e^{-\rho X_\alpha}.
\]
This Gaussian lattice sum decays exponentially at infinity. Poisson summation in the finitely many residue classes gives rapid decay at zero as well, since the degree-one polynomial makes the dual zero-frequency integral zero. Let $L_j^{(1)}$ be the associated Laguerre polynomial and set
\[
 B_j=\sum_{\alpha\in\Lambda}w(\alpha)\alpha L_j^{(1)}(2X_\alpha)e^{-X_\alpha}.
\]
For $t=(\rho-1)/(\rho+1)$ the generating function
$\sum_jL_j^{(1)}(x)t^j=(1-t)^{-2}\exp(-xt/(1-t))$ gives $\Theta(\rho)=(1-t)^2\sum_jB_jt^j$. Consequently for every integer $h\ge0$,
\[
 \int_0^\infty\Theta(\rho)(\log\rho)^h\,d\rho
 =2\lim_{v\uparrow1}\sum_{j\ge0}v^jB_j
 \int_{-1}^1 t^j\left(\log\frac{1+t}{1-t}\right)^h dt .
\]
For $0<v<1$ the power series converges uniformly on $[-v,v]$. Its sum is $\Theta((1+vt)/(1-vt))/(1-vt)^2$, uniformly bounded by endpoint decay. The remaining logarithmic power is integrable, so dominated convergence gives the formula. After assembling the ideal-class sums and conductor powers in the Hecke Mellin formula, these integrals give ordinary complex derivatives at the centre. This uses the radial coefficients $B_j$; Nomoto's scalar recurrence in the angular weight supplies the central test \eqref{test}. On the branch $p\mid F_p$, the exact algebraic and analytic ranks can still require further information.

\begin{thebibliography}{9}\raggedright
\bibitem{nom} K.~Nomoto, \emph{The rank of a CM elliptic curve and a recurrence formula}, J. Number Theory (2022). \href{https://arxiv.org/abs/2103.08947}{Author's paper}.
\bibitem{thesis} K.~Nomoto, \emph{The $p$-adic valuations of the critical values of $L$-functions associated to elliptic curves}, Kyushu University dissertation, Part II, Propositions 5.6 and 5.8, Remark 5.9 and Theorem 7.6. \href{https://catalog.lib.kyushu-u.ac.jp/opac_download_md/6787428/math0277.pdf}{University repository}.
\bibitem{mr} B.~Mazur and K.~Rubin, \emph{Ranks of twists of elliptic curves and Hilbert's tenth problem}, Invent. Math. \textbf{181} (2010), 541--575. \href{https://arxiv.org/abs/0904.3709}{Authors' paper}.
\bibitem{sil} J.~H.~Silverman, \emph{The Arithmetic of Elliptic Curves}, 2nd ed., Springer (2009), X.4. \href{https://doi.org/10.1007/978-0-387-09494-6}{Publisher}.
\bibitem{ps} B.~Poonen and M.~Stoll, \emph{The Cassels--Tate pairing on polarized abelian varieties}, Ann. of Math. \textbf{150} (1999), 1109--1149. \href{https://math.mit.edu/~poonen/papers/sha.pdf}{Authors' paper}.
\bibitem{bt} A.~A.~Burungale and Y.~Tian, \emph{A rank zero $p$-converse to a theorem of Gross--Zagier, Kolyvagin and Rubin}, Ann. of Math. \textbf{203} (2026), 1--13. \href{https://doi.org/10.4007/annals.2026.203.1.1}{Publisher}.
\bibitem{creutz} B.~Creutz, \emph{There are no transcendental Brauer--Manin obstructions on abelian varieties}, Int. Math. Res. Not. (2020), 2684--2697. \href{https://doi.org/10.1093/imrn/rny098}{Publisher}.
\bibitem{pfBF} A.~Burungale and M.~Flach,
\emph{The conjecture of Birch and Swinnerton-Dyer for certain elliptic
curves with complex multiplication}, \href{https://arxiv.org/abs/2206.09874}{Corollary~2}.
\bibitem{pfSmith} A.~Smith, \emph{The distribution of $\ell^\infty$-Selmer
groups in degree $\ell$ twist families I},
\href{https://arxiv.org/abs/2207.05674}{Theorem~1.5}.
\bibitem{pfMorgan} A.~Morgan, \emph{Hasse principle for Kummer varieties
in the case of generic 2-torsion},
\href{https://doi.org/10.1112/plms.70066}{Sections~2 and 5.2}.
\bibitem{pfBS} A.~Burungale and C.~Skinner,
\emph{A $p$-converse theorem for CM elliptic curves}, Appendix~A to
Alp\"oge--Bhargava--Shnidman, \href{https://people.math.harvard.edu/~alpoge/papers/integers%20expressible%20as%20the%20sum%20of%20two%20rational%20cubes.pdf}{Theorem~A.1 and Remark~A.5}.
\end{thebibliography}
\end{document}